% Appendix B.1, input by main.tex.

\FloatBarrier
% ============================================================
\section{Primitive Composition}\label{sec:composition}
% ============================================================

\subsection{Prompt structure overview}\label{sec:composition-overview}

In primitive composition, the~\gls{ra} prompts the~\gls{llm} to compose each accepted recovery event from the
primitives of the resource that performs it. Regardless of the generation level, events are
composed one at a time, in the order of the accepted outline. In each call, the~\gls{llm} requests facts, returns a program consisting of an ordered
list of primitive steps, or reports that the event cannot be composed from the primitives.
The~\gls{llm} supplies numeric parameters as references to facts provided by the~\gls{ra}, which resolves
those references. The~\gls{ra} then checks whether the program reaches the event's intended end state
and whether each step can be carried out in the cell. If the program is refused, the next prompt
includes the validation finding so that the~\gls{llm} can revise the program. This appendix describes
the prompt for one event. Table~\ref{tab:composition-sections} lists its sections in the order in
which they appear.

% A table that may continue on the next page, so the text runs on without gaps.
{\footnotesize\setlength{\tabcolsep}{5pt}
\begin{xltabular}{\linewidth}{@{}>{\raggedright\arraybackslash}p{4.4cm} >{\raggedright\arraybackslash}p{2.4cm} Y@{}}
\toprule
Section & In the prompt & Description \\
\midrule
\endfirsthead
\toprule
Section & In the prompt & Description (continued) \\
\midrule
\endhead
\bottomrule
\endfoot
\bottomrule
\caption{The sections of the primitive composition prompt, in the order in which they appear.}\label{tab:composition-sections}\\
\endlastfoot
Task and Role & every call & Asks the~\gls{llm} to compose only the active recovery event, using its resource's primitives, and to request missing facts rather than guess them. \\
Primitives & every call & The names of the primitives of the resource that performs the event. \\
Knowledge Base & every call & Available lookups, such as a primitive's contract, the resource's model, a part's observed pose and target, area bounds and the geometry of the cell. \\
Nominal functions & every call & The names of the resource's nominal functions, whose primitive steps can be requested. \\
Requested facts & after a lookup & The facts requested by the~\gls{llm} while composing this event, such as the resource's model and primitive contracts. \\
Validation feedback & after a refusal & The finding from checking the previous program, identifying the step or end-state value that needs revision. \\
Active event & every call & The recovery event as accepted, including its resource and part, its starting state, its intended end state and the rationale provided by the~\gls{llm}. \\
Predefined configurations & every call & Configurations declared by the resource, such as home or transport, that a step can reference by name. The list is empty for a resource that declares none, such as a machine. \\
Composition Rules & every call & Requirements for the reply and an accepted program: use only the listed primitives, express every number as a reference to a requested fact, and establish every intended end-state value through declared effects or what the checks find. \\
System message & every call & The required reply format: a short plan accompanied by a request for facts, a program of primitive steps, or a report that the event cannot be composed from the primitives. \\
\end{xltabular}}
